\chapter{Conclusion}

During this project, we implemented and went into the details of three different recommender systems for Spotify songs that are based on two different paradigms:  We started with content-based filtering using a nearest neighbors model (cf. Chapter \ref{chap:nn}) and followed with collaborative filtering using on the one hand a Logistic Matrix Factorization (cf. Chapter \ref{chap:lmf}) and on the other hand Variational Autoencoders (cf. Chapter \ref{chap:vae}). 

With our objective being to be able to measure and compare the effectiveness of content-based filtering and collaborative filtering in recommender systems while dealing with a small user base. We implemented a metric: the Vector Set Divergence outlined in Section \ref{sec:vsd} to be able to quantitatively compare the different outputs of all our models. Additionally, we also implemented a data generator as discussed in Section \ref{sec:data_generation}, to address the size limitation of the dataset we retrieved from Spotify, and have the flexibility to control the quantity of data and improve our experiment.


Our results as detailed in Chapter \ref{chap:results} show that, based on the Vector Set Divergence metric, the collaborative-filtering paradigm despite the constraint of a small user base achieves in general better scores, with the VAE and LMF models scoring very similar scores, than the content-based filtering. It is also worth noting that the Nearest Neighbors model tend to have very slow training times compared to both other models.

In future work, investigating hybrid approaches that combine both content-based filtering and collaborative-filtering may increase the accuracy and inference time. Adding a human-based metric could also enhance the overall accuracy of the models and ensure the models serve their purpose.